% !PRINTED 533
\weberpage{533}

\begin{center}
\S.\,162.

\gesperrt{Abelian Equations.}
\end{center}

The degree of a transitive permutation group of $m$ figures
is always divisible by $m$, and therefore never smaller than $m$.
For if $P$ is such a group and $Q_0$ is the totality of the permutations
of $P$ that leave the figure 0 unchanged, then $Q_0$ is also
a group, and therefore a divisor of $P$. Now, on account of the
assumed transitivity, one can find in $P$ a system of
permutations $\pi_1, \pi_2 \ldots \pi_{m-1}$ that carry 0 into 1, 0 into 2, \ldots
% [sic] authority prints Q\pi_1 here, omitting the subscript 0 expected from Q_0.
0 into $m-1$, and then $Q\pi_1$ is the system of all those
permutations of $P$ that change 0 into 1. Accordingly

\weq{1}{P = Q_0 + Q_0\pi_1 + Q_0\pi_2 + \cdots + Q_0\pi_{m-1},}

and so the degree of $P$ is equal to the product of $m$ and the degree
of $Q_0$.

\gesperrt{The Galois group of an irreducible equation is therefore never of lower degree than the equation itself.}

A normal equation we have earlier defined by the stipulation
that it be irreducible and that each of its
roots should be rationally expressible in terms of every other
(\S.\,145). From this it follows that the group of a normal equation
must reduce to the identical group when
one root is adjoined; and conversely, an equation
whose group is so constituted is, provided it is at the same time irreducible,
always a normal equation.

Now $P$ reduces, by adjunction of the root $\alpha_0$, to
$Q_0$, by $\alpha_1$ to $\pi_1^{-1}Q_0\pi_1$, etc., and so if $P$ is to be the group
of a normal equation, it is necessary and sufficient
that $Q_0$ be the unit group, and hence that the group be exhausted by
% [sic] authority prints an unindexed \pi before \pi_1,...,\pi_{m-1}; none is introduced here, and context calls for the identity representative.
$\pi, \pi_1, \pi_2, \ldots \pi_{m-1}$. We therefore have:

\gesperrt{In order that an irreducible equation be a normal equation, it is necessary and sufficient that the degree of the group agree with the degree of the equation.}

We first consider here the special class of equations that includes
the cyclotomic equations first solved by \gesperrt{Gauss},
%CONT>>
